\documentclass[11pt,a4paper]{article}

\usepackage[T1]{fontenc}
\usepackage[utf8]{inputenc}
\usepackage[portuguese]{babel}
\usepackage{lmodern}
\usepackage{amsmath,amssymb}
\usepackage{booktabs}
\usepackage{enumitem}
\usepackage[a4paper,margin=2.4cm]{geometry}
\usepackage[hidelinks]{hyperref}
\usepackage{microtype}

\setlength{\parindent}{0pt}
\setlength{\parskip}{0.55em}
\setlength{\emergencystretch}{4em}
\sloppy
\setlist[itemize]{leftmargin=1.6em,itemsep=0.3em,topsep=0.3em}
\setlist[enumerate]{leftmargin=1.8em,itemsep=0.45em,topsep=0.3em}

\newcommand{\localizacao}{\textbf{Localização.}\ }
\newcommand{\problema}{\textbf{Problema.}\ }
\newcommand{\impacto}{\textbf{Por que importa.}\ }
\newcommand{\correcao}{\textbf{Correção sugerida.}\ }

\title{Parecer técnico sobre\\
\emph{Optimal Control of Static Var Compensators for Power Quality Enhancement in Three-Phase Distribution Systems}}
\author{}
\date{16 de julho de 2026}

\begin{document}
\maketitle

\section*{Avaliação geral}

O manuscrito trata de um problema relevante e tenta integrar fluxo de potência
harmônico trifásico, controle independente dos ramos delta de SVCs FC--TCR e
restrições de tensão, THD e desequilíbrio. A organização conceitual é
promissora, mas a versão atual requer \textbf{revisão maior} antes que os
resultados possam sustentar as conclusões. Há inconsistências quantitativas nos
parâmetros dos SVCs, erros ou ambiguidades em fórmulas usadas pela implementação,
conflitos na definição do sistema-teste e falta de comparações que isolem o
benefício do controle por ramo.

Foram verificados explicitamente os sinais, as dimensões, as convenções de
ramo/fase, a decomposição em componentes simétricas, o truncamento harmônico,
as janelas de condução e as fórmulas voltadas à implementação. Também foram
comparados resumo, metodologia, tabelas, apêndice e conclusão. Os 22 arquivos
PDF chamados por \verb|\includegraphics| não estão presentes no conjunto de
arquivos fornecido; por isso, as curvas e seus valores não puderam ser
inspecionados diretamente. A tabela numérica e o texto, contudo, já revelam os
problemas abaixo.

\section*{Achados técnicos prioritários}

\subsection*{1. Potências e parâmetros dos SVCs não são dimensionalmente compatíveis}

\localizacao Subseção ``Fixed-Capacitor SVC Representation and Rated Power'',
Eq.~\texttt{eq:svc\_rating}; Tabela~\texttt{tab:svc\_parameters};
Eq.~\texttt{eq:svc\_total\_rating}; Apêndice,
Tabela~\texttt{tab:appendix\_svc\_data}.

\problema \textbf{Erro definido.} Com $C=2652{,}6\,\mu\mathrm{F}$ e
$f=60\,\mathrm{Hz}$, obtém-se
$B_C=2\pi fC\simeq1{,}000\,\mathrm{S}$ e $X_C\simeq1\,\Omega$, como a
própria tabela declara. Para três ramos delta a $V_{LL}=12{,}66\,\mathrm{kV}$,
a potência capacitiva no bloqueio completo do TCR seria
\[
  Q_C=3V_{LL}^2B_C\simeq480{,}83\,\mathrm{MVAr},
\]
e não $3{,}2108$ ou $3{,}0651\,\mathrm{MVAr}$. A discrepância é superior a
duas ordens de grandeza e não pode ser explicada pela pequena diferença entre
as tensões das barras. Além disso, a soma de
$\lvert B_{\mathrm{SVC},j,\ell}\rvert$ em
\texttt{eq:svc\_rating} produz uma grandeza sempre não negativa, enquanto a
tabela informa um limite indutivo negativo. Ela também deixa de representar a
potência reativa líquida quando ramos têm sinais diferentes.

\impacto O dimensionamento dos compensadores determina as tensões, os ângulos
ótimos, a THD e a aparente melhoria de carregabilidade. Se os dados usados na
simulação tiverem a mesma inconsistência, todos os resultados precisam ser
recalculados.

\correcao Conferir a unidade de $C_j$, a tensão efetivamente aplicada aos
ramos e as bases por unidade. Para potência líquida, usar uma soma com sinal,
por exemplo $Q_j=V_{LL}^2\sum_\ell
B_{\mathrm{SVC},j,\ell}$, com uma convenção capacitiva/indutiva declarada.
Para ``capacidade instalada'', reportar separadamente as potências nominais do
FC e do TCR, ou $\max_\alpha |Q_j(\alpha)|$, sem transformar a potência líquida
em soma de módulos. Recalcular a Tabela~\texttt{tab:svc\_parameters} a partir
dos mesmos dados efetivamente usados pelo fluxo harmônico.

\subsection*{2. O modelo de carga perde o sinal da geração distribuída}

\localizacao Eqs.~\texttt{eq:pq\_current\_injection}--
\texttt{eq:pq\_current\_imag} e Tabela~\texttt{tab:appendix\_load\_data}.

\problema \textbf{Erro definido.} As componentes retangulares substituem
$P_{i,\phi}$ e $Q_{i,\phi}$ por $|P_{i,\phi}|$ e $|Q_{i,\phi}|$. Entretanto,
o apêndice contém potências ativas negativas após a subtração de geração, por
exemplo $90-200$ na barra 18. O módulo converte essa injeção líquida em uma
carga positiva. Há ainda uma ambiguidade de sinal: a equação chama $I$ de
corrente injetada, mas usa potência de demanda positiva sem o sinal negativo
normalmente necessário nessa convenção.

\impacto O erro altera o fluxo fundamental, as perdas, a queda de tensão, a
sequência negativa e o ponto ótimo. Também invalida a interpretação dos casos
com geração distribuída.

\correcao Manter $P$ e $Q$ com sinal nas expressões retangulares e declarar
uma única convenção. Se $S^{\mathrm{dem}}=P+jQ$ for demanda positiva, usar
$I^{\mathrm{inj}}=-(S^{\mathrm{dem}}/V)^*$; alternativamente, armazenar cargas
como injeções negativas. Atualizar de modo coerente o resíduo e o Jacobiano e
adicionar um teste unitário com uma barra de carga e outra de geração.

\subsection*{3. A transformação harmônica de sequências é singular para ordens incluídas no estudo}

\localizacao Eq.~\texttt{eq:harmonic\_sequence\_decomposition} e conjunto
$\mathcal H=\{1,3,5,7,11,13\}$ no apêndice.

\problema \textbf{Erro definido.} A matriz usa $a^h$ e $a^{2h}$. Para
$h=3$, todas as três linhas se tornam $[1\;1\;1]/3$; portanto, a matriz tem
posto um e não decompõe a corrente em sequências zero, positiva e negativa.
Para outras ordens, a dependência em $h$ também troca os rótulos das sequências
em relação à transformação de Fortescue usual, sem que uma convenção
alternativa seja definida.

\impacto A ordem 3 faz parte da simulação, e a decomposição é usada para
interpretar os efeitos do controle assimétrico. Uma classificação incorreta
pode ocultar componentes negativas ou atribuí-las à sequência errada.

\correcao Aplicar a matriz de Fortescue não singular e independente de $h$ a
cada vetor de fasores harmônicos,
\[
 \frac{1}{3}
 \begin{bmatrix}1&1&1\\1&a&a^2\\1&a^2&a\end{bmatrix},
\]
e depois interpretar a sequência característica de cada ordem. Se o código
usa outra convenção, derivá-la, demonstrar que é inversível para todo
$h\in\mathcal H$ e alinhar os rótulos $+$ e $-$.

\subsection*{4. A THD calculada não corresponde claramente ao conjunto harmônico nem basta para alegar conformidade}

\localizacao Eq.~\texttt{eq:thd}, subseções ``Numerical Settings'' e
``Three-Phase Unbalanced Loading Condition'', e apêndice.

\problema \textbf{Erro de notação e alegação não sustentada.} Primeiro, $H$
é definido como o número de elementos de $\mathcal H$ e depois como a maior
ordem, enquanto a soma $\sum_{h=2}^{H}$ sugere ordens consecutivas que não são
modeladas. Segundo, no caso assimétrico o texto reconhece que o cancelamento
ideal de triplas deixa de ser garantido, mas $\mathcal H$ omite a 9ª harmônica.
O conjunto também termina na 13ª ordem e não é apresentada análise de
convergência espectral. Assim, o valor reportado é, no máximo, uma THD parcial
no conjunto selecionado.

\impacto O otimizador pode respeitar uma restrição subestimada e o texto pode
declarar conformidade de 8\% mesmo com componentes omitidas relevantes.

\correcao Escrever
$\mathrm{THD}_{V,i,\phi}=100
\sqrt{\sum_{h\in\mathcal H\setminus\{1\}}|V_{i,\phi}^{(h)}|^2}/
|V_{i,\phi}^{(1)}|$. Incluir ao menos as ordens necessárias ao padrão adotado,
inclusive triplas que não se cancelam sob controle desigual, e demonstrar
sensibilidade ao aumento de $h_{\max}$. Até isso ser feito, substituir
``conformidade'' por ``satisfação do limite no espectro truncado modelado''.

\subsection*{5. Fasor, magnitude, fase e ramo compartilham símbolos incompatíveis}

\localizacao Eqs.~\texttt{eq:complex\_voltage\_vector},
\texttt{eq:optproblem}, \texttt{eq:obj\_dev}, \texttt{eq:obj\_imb} e
\texttt{eq:thd}; Eqs.~\texttt{eq:tcr\_current\_time\_solution}--
\texttt{eq:tcr\_imag\_fourier}; discussão e legendas dos ângulos no caso
desequilibrado.

\problema \textbf{Ambiguidade de implementação.} $V_{i,\phi}^{(h)}$ aparece
primeiro como fasor complexo e depois como magnitude real em desigualdades e
diferenças de fase, sem barras de módulo. As desigualdades
$V_{\min}\le V_{i,\phi}^{(1)}\le V_{\max}$ não são definidas para números
complexos. Em paralelo, a corrente do ramo $\ell$ passa a ser indexada por
$\phi$ nas fórmulas do TCR, e os resultados chamam $\alpha_{j,ab}$,
$\alpha_{j,bc}$ e $\alpha_{j,ca}$ de ângulos ``entre fases'' ou ``por fase''.
O descritor inicial é \emph{ângulo do ramo delta}; o descritor posterior é
\emph{ângulo de fase}, que não é equivalente.

\impacto Essas mudanças podem levar diretamente à implementação do objetivo
com fasores em vez de magnitudes ou à aplicação de controles $a,b,c$ no lugar
dos ramos $ab,bc,ca$.

\correcao Reservar, por exemplo, $\underline V_{i,\phi}^{(h)}$ para o fasor,
$v_{i,\phi}^{(h)}=|\underline V_{i,\phi}^{(h)}|$ para a magnitude e
$\theta_{i,\phi}^{(h)}$ para o ângulo. Usar $v$ nas restrições, em
$J_{\mathrm{dev}}$, $J_{\mathrm{imb}}$ e na THD. Manter $\ell$ em todas as
equações do ramo TCR e nomear as três curvas de controle como $ab$, $bc$ e
$ca$. Declarar também se $|V_2|$ nos resultados é exatamente $|V^-|$ da
metodologia.

\subsection*{6. A ordenação do vetor de estado não coincide com a matriz retangular exibida}

\localizacao Eqs.~\texttt{eq:mict\_bus\_state},
\texttt{eq:mict\_state} e \texttt{eq:real\_ybus\_h}.

\problema \textbf{Erro definido, salvo se houver uma permutação omitida.} O
estado por barra é intercalado como
$[V_{a,r},V_{b,r},V_{c,r},V_{a,m},V_{b,m},V_{c,m}]$ e depois empilhado por
barra. Já a matriz
$\begin{bmatrix}\Re Y&-\Im Y\\\Im Y&\Re Y\end{bmatrix}$ pressupõe o arranjo
global $[\mathbf V_r;\mathbf V_m]$. Para mais de uma barra, essas ordenações
não são iguais.

\impacto Sem uma matriz de permutação, os produtos matriz--vetor e o Jacobiano
associam tensões e equações a posições erradas.

\correcao Escolher uma única ordenação global ou introduzir explicitamente a
matriz de permutação $P$ e usar $P^T\mathcal YP$. Incluir no texto um exemplo
de duas barras e um teste que compare a forma complexa com a forma real.

\subsection*{7. As janelas de condução não são coerentes com o modelo $R$--$L$ harmônico}

\localizacao Eqs.~\texttt{eq:tcr\_firing\_time\_positive}--
\texttt{eq:tcr\_extinction\_time\_negative},
\texttt{eq:tcr\_rl\_differential}--
\texttt{eq:tcr\_initial\_constant}, e descrição das janelas no apêndice.

\problema \textbf{Provável erro de implementação.} As extinções fixas em
$2\pi-\alpha$ e $3\pi-\alpha$ valem para a aproximação ideal, senoidal e
essencialmente sem perdas. O próprio modelo, porém, usa $R>0$ e uma tensão com
várias harmônicas; nessas condições o instante natural de corrente nula deve
ser resolvido, não imposto pela simetria ideal. Além disso, o apêndice declara
$T_{\mathrm{on},2}=\alpha_j$, enquanto a metodologia define o segundo disparo
em $\pi+\alpha_{j,\ell}$. Não se informa se há um deslocamento implícito. A
reconstrução temporal também não esclarece se os fasores são RMS ou de pico,
isto é, se falta um fator $\sqrt 2$.

\impacto Os coeficientes de Fourier, a susceptância fundamental e a THD
dependem diretamente das janelas e da escala temporal. Isso pode explicar
parte da inconsistência numérica dos SVCs.

\correcao Ou assumir formalmente $R=0$ e tensão senoidal para determinar as
janelas, usando o modelo harmônico apenas depois, ou localizar
$t_f$ pela primeira raiz de $i(t)$ após o disparo para a forma de onda
reconstruída. Harmonizar o segundo disparo com $\pi+\alpha$ e declarar a
convenção RMS/pico em todas as integrais.

\subsection*{8. Sistema-teste, numeração e topologia são internamente contraditórios}

\localizacao Subseção ``System Description and Test Network'',
Fig.~\texttt{fig:ieee33\_singleline} e Tabelas
\texttt{tab:appendix\_line\_data}--\texttt{tab:appendix\_transformer\_data}.

\problema \textbf{Erro definido e lacuna de reprodutibilidade.} O texto e o
desenho colocam os SVCs controlados nas barras 17 e 32, mas a legenda da figura
declara barras 16 e 31 ``under the shifted numbering scheme''. Além disso, a
tabela de linhas já contém 5--6 e 7--8, enquanto a tabela de transformadores
repete essas ligações e acrescenta 4--9, 7--9 e 4--7. Se todos esses elementos
forem ativos, a rede não coincide com o alimentador radial desenhado. Os códigos
de conexão ``11'', unidades, relações de transformação e bases dos parâmetros
também não são definidos. A afirmação de que capacitâncias shunt são
``derivadas'' de indutâncias mútuas não fornece uma relação dimensional ou uma
equação reproduzível.

\impacto Não é possível reconstruir a matriz $Y^{(h)}$ nem saber quais barras
recebem controle. Comparações com o benchmark citado ficam indeterminadas.

\correcao Fornecer uma tabela única de mapeamento entre numeração original e
numeração deslocada; corrigir texto, desenho, legenda e arquivos de dados para
o mesmo par de barras. Explicar se a tabela de transformadores substitui linhas
ou adiciona ramos, mostrar esses elementos no diagrama e definir conexão,
relação, unidade e base. Publicar a equação usada para os termos mútuos e shunt.

\subsection*{9. A tabela de resultados mistura os limites de VUF}

\localizacao Subseção ``Three-Phase Unbalanced Loading Condition'' e
Tabela~\texttt{tab:unbalanced\_performance\_summary}.

\problema \textbf{Erro definido.} O caso desequilibrado adota explicitamente
$\mathrm{VUF}_{\max}=1\%$. A tabela informa VUF máximo sem SVC de
$1{,}5876\%$ na barra 32, mas marca ``VUF limit, no SVC'' como ``Yes''. Esse
valor viola o limite de 1\%; a célula só seria ``Yes'' sob o limite geral de
2\%, que não é o critério declarado para esse caso.

\impacto A linha de conformidade resume a principal métrica do artigo e, como
está, induz a uma conclusão errada.

\correcao Alterar a célula para ``No'' e indicar no cabeçalho da tabela que o
limite do caso é 1\%. Não misturar o limite geral de 2\% com o alvo do ensaio de
estresse. Regerar automaticamente as linhas ``Yes/No'' a partir dos valores e
do limite usado na simulação.

\subsection*{10. Os cenários não isolam a contribuição do controle por ramo e não demonstram carregabilidade}

\localizacao Resumo, resultados equilibrados e desequilibrados e conclusão.

\problema \textbf{Alegação não sustentada.} O manuscrito afirma aumento da
capacidade de carregamento, mas não fornece $\lambda_{\max}$, critério de
colapso/inviabilidade ou comparação numérica equilibrada sem controle. No caso
desequilibrado, compara-se ``sem SVC'' com ``SVC otimizado'', sem um caso de
ângulo comum. Portanto, a melhora pode decorrer da mera instalação de potência
reativa, e não da independência entre $ab$, $bc$ e $ca$. Como SQP é um método
local e não há análise multi-inicialização, ``optimal'' também deve ser
entendido como ótimo local, não global.

\impacto As duas contribuições centrais---maior carregabilidade e vantagem do
controle por ramo---não são identificadas pelo desenho experimental atual.

\correcao Incluir, com os mesmos equipamentos e limites: (i) sem SVC;
(ii) SVC com ângulo fixo; (iii) otimização com um ângulo comum por SVC; e
(iv) otimização independente por ramo. Reportar $\lambda_{\max}$, tensão
mínima, VUF, THD, potência reativa e ângulos, além do critério que encerra a
curva de carga. Qualificar a solução como local ou apresentar evidência de
robustez a diferentes pontos iniciais.

\subsection*{11. Dados e algoritmo ainda não permitem reprodução independente}

\localizacao Subseções ``Phase-Dependent Load Curve'', ``Sequential
Quadratic Programming Solver'' e ``Solution Algorithm''; definição do caso
desequilibrado; apêndice.

\problema \textbf{Ambiguidade de implementação.} Os 24 valores dos perfis
$\lambda_{i,\phi,k}$ não são tabulados; a figura ausente não substituiria os
dados exatos. Não se especificam o software SQP, cálculo dos gradientes,
ângulos iniciais, tolerâncias de viabilidade separadas, tratamento de falhas ou
estatísticas de convergência. O algoritmo conserva os ângulos anteriores quando
o SQP falha, mas não verifica se esses ângulos continuam viáveis no novo nível
de carga. Também não fica claro se ``without SVC'' remove apenas os SVCs das
barras 17 e 32 ou mantém o SVC fixo da barra 11; a afirmação de que os dois SVCs
controlados são as únicas fontes harmônicas conflita com a descrição anterior
de três dispositivos.

\impacto Não é possível repetir o experimento nem confirmar que todos os
pontos apresentados satisfazem as restrições.

\correcao Fornecer os perfis horários em tabela/arquivo, a função completa de
$B_{\mathrm{TCR}}(\alpha)$, os dados do caso-base e a definição exata de cada
cenário. Documentar solver, derivadas, inicialização e teste de viabilidade.
Quando o SQP falhar, marcar o ponto como não resolvido ou executar uma fase de
restauração; não armazená-lo como solução válida sem nova checagem. Reportar
número de falhas, iterações e tempo computacional.

\section*{Questões normativas e factuais a verificar externamente}

\begin{itemize}
  \item \textbf{IEC 61000-2-2:} o texto atribui o limite de 8\% a sistemas
  públicos de média tensão, mas o próprio título bibliográfico citado diz
  ``Public Low-Voltage Power Supply Systems''. Trata-se de uma contradição
  interna definida. Confirmar na edição oficial vigente se o caso de
  $12{,}66\,\mathrm{kV}$ deve citar a IEC 61000-2-12 ou outro documento e quais
  ordens/condições integram a avaliação de THD.
  \item \textbf{EN 50160:} confirmar na edição oficial vigente se 2\% é limite
  absoluto aplicável a cada ponto ou valor estatístico sob uma janela de
  medição definida. O manuscrito deve informar edição, agregação temporal e
  critério de conformidade; o ensaio interno de 1\% deve ser distinguido da
  exigência normativa.
  \item \textbf{Benchmark IEEE de 33 barras:} conferir no material primário de
  \texttt{dolatabadi2021} a numeração deslocada, os dados por fase e o papel dos
  cinco transformadores. A versão atual não permite decidir se as tabelas são
  adaptação intencional ou mistura de dois modelos.
\end{itemize}

\section*{Achados editoriais significativos}

\begin{itemize}
  \item \localizacao Resumo. \problema O trecho exato
  \texttt{formulated sed on} é uma frase corrompida. \impacto O erro aparece
  na apresentação principal do método. \correcao Substituir por
  \texttt{formulated based on}.
  \item \localizacao Metodologia e resultados. \problema Alternam-se
  ``branch-wise'', ``phase-wise'', ``phase firing angles'' e ``effective firing
  angles'' para variáveis que pertencem aos ramos delta. \impacto A variação
  reforça a ambiguidade matemática descrita no Achado~5. \correcao Usar
  sistematicamente ``delta-branch firing angles'' e identificar $ab$, $bc$ e
  $ca$ nas legendas.
  \item \localizacao Resultados desequilibrados. \problema ``without-SVC'' não
  define quais dos três compensadores são removidos. \impacto O leitor não sabe
  qual é o caso-base. \correcao Nomear os cenários de forma explícita, por
  exemplo ``all SVCs disconnected'', ``fixed bus-11 SVC only'' e ``controlled
  bus-17/bus-32 SVCs enabled''.
\end{itemize}

\subsection*{Varredura final de prova}

O script fornecido \texttt{scripts/proofing\_scan.py} foi executado sobre
\texttt{main.tex} com limite de 80 ocorrências e não encontrou casos dos
padrões mecânicos previstos (pontuação duplicada, capitalização após
ponto-e-vírgula, molduras ``into\ldots into'', nomes de produto ou
\texttt{arctan(x/y)}). A inspeção manual encontrou os seguintes itens de alta
confiança:

\begin{itemize}
  \item \localizacao Resumo. \problema \texttt{formulated sed on}.
  \correcao \texttt{formulated based on}.
  \item \localizacao Figura~\texttt{fig:ieee33\_singleline}.
  \problema O rótulo interno ``A-C System'' usa uma grafia incomum para
  corrente alternada. \correcao Usar ``AC system''.
  \item \localizacao Bibliografia. \problema Há 18 itens, mas o ambiente é
  aberto com \verb|\begin{thebibliography}{1}|; além disso,
  \verb|\bibliographystyle{IEEEtran}| está dentro de uma bibliografia manual e
  não tem efeito útil. \correcao Usar \verb|{99}| e remover o comando de estilo,
  ou migrar integralmente para BibTeX.
  \item \localizacao Bibliografia. \problema As entradas
  \texttt{Hu2023HarmonicSources} e \texttt{ANEELProdist8} não são citadas.
  \correcao Citá-las em afirmações que realmente sustentem ou removê-las.
\end{itemize}

Não foram observadas pontuações duplicadas ou aspas de títulos manifestamente
malformadas nas referências.

\section*{Ordem recomendada das correções}

\begin{enumerate}
  \item Reconciliar unidades, bases e potências dos SVCs e confirmar os dados
  realmente usados na simulação.
  \item Corrigir o sinal das cargas/geração, a decomposição de sequências, a
  soma de THD, a ordenação retangular e as janelas do TCR; validar cada bloco
  com testes simples antes de rerodar a otimização.
  \item Tornar inequívocos a topologia, a numeração das barras, os
  transformadores, os perfis horários e a definição de cada cenário.
  \item Rerodar os quatro casos comparativos, corrigir a tabela de VUF e
  reportar valores numéricos de carregabilidade, viabilidade e convergência.
  \item Só então revisar resumo e conclusão, limitando as alegações ao que as
  comparações e o espectro harmônico efetivamente demonstram, e concluir a
  revisão editorial/bibliográfica.
\end{enumerate}

\end{document}